\subsectionaddtolist{الف}
\[
X(s) = \frac{1}{(s + 1)(s + 2)}
\]
این تابع دارای دو قطب در نقاط $s = -1$ و $s = -2$ است. ابتدا آن را ساده می‌کنیم:
\[
X(s) = \frac{1}{s + 1} - \frac{1}{s + 2}
\]

با توجه به مکان قطب‌ها، سه ناحیه همگرایی (ROC) ممکن وجود دارد که به ازای هر کدام، سیگنال زمانی $x(t)$ را محاسبه می‌کنیم:

\begin{enumerate}
    \item $\text{Re}(s) > -1:$

    هر دو قطب به صورت راست‌بر (بر حسب تابع پله $u(t)$) ظاهر می‌شوند:
    \[
    x_1(t) = e^{-t} u(t) - e^{-2t} u(t) = \left( e^{-t} - e^{-2t} \right) u(t)
    \]

    \item $\text{Re}(s) < -2:$

    هر دو قطب به صورت چپ‌بر (بر حسب تابع پله $-u(-t)$) ظاهر می‌شوند:
    \[
    x_2(t) = -e^{-t} u(-t) - \left( -e^{-2t} u(-t) \right) = \left( e^{-2t} - e^{-t} \right) u(-t)
    \]

    \item $-2 < \text{Re}(s) < -1:$

    قطب $s = -1$ چپ‌بر و قطب $s = -2$ راست‌بر خواهد بود:
    \[
    x_3(t) = -e^{-t} u(-t) - e^{-2t} u(t)
    \]
\end{enumerate}


\subsectionaddtolist{ب}

شرط همگرایی تبدیل لاپلاس، محدود بودن انتگرال قدرمطلق سیگنال وزن‌دهی شده با نمایی حقیقی است:
\[
I = \int_{-\infty}^{+\infty} |x(t)| e^{-\sigma t} dt < \infty
\]

این انتگرال را برای هر سه سیگنال به دست آمده بررسی می‌کنیم:

\begin{enumerate}
    \item $x_1(t) = (e^{-t} - e^{-2t})u(t)$:

    چون برای $t>0$ همواره $e^{-t} \ge e^{-2t}$ است، قدرمطلق خود عبارت درون پرانتز است:
    \[
    I_1 = \int_{0}^{+\infty} (e^{-t} - e^{-2t}) e^{-\sigma t} dt = \int_{0}^{+\infty} e^{-(\sigma + 1)t} dt - \int_{0}^{+\infty} e^{-(\sigma + 2)t} dt
    \]
    برای همگرایی این انتگرال در بی‌نهایت، باید توان‌های نمایی منفی باشند:
    \[
    \begin{cases} \sigma + 1 > 0 \implies \sigma > -1 \\ \sigma + 2 > 0 \implies \sigma > -2 \end{cases} \implies \text{ :اشتراک} \sigma > -1
    \]

    \item $x_2(t) = (e^{-2t} - e^{-t})u(-t)$:

    برای زمان‌های منفی $t<0$ همواره $e^{-2t} \ge e^{-t}$ است:
    \[
    I_2 = \int_{-\infty}^{0} (e^{-2t} - e^{-t}) e^{-\sigma t} dt = \int_{-\infty}^{0} e^{-(\sigma + 2)t} dt - \int_{-\infty}^{0} e^{-(\sigma + 1)t} dt
    \]
    برای همگرایی این انتگرال در منفی بی‌نهایت، باید ضریب $t$ مثبت باشد تا در نهایت منفی در منفی مثبت شده و مایه واگرایی نشود:
    \[
    \begin{cases} \sigma + 2 < 0 \implies \sigma < -2 \\ \sigma + 1 < 0 \implies \sigma < -1 \end{cases} \implies \text{ :اشتراک} \sigma < -2
    \]

    \item $x_3(t) = -e^{-t}u(-t) - e^{-2t}u(t)$:

    انتگرال را به دو بازه مثبت و منفی تجزیه می‌کنیم:
    \[
    I_3 = \int_{-\infty}^{0} |-e^{-t}| e^{-\sigma t} dt + \int_{0}^{+\infty} |-e^{-2t}| e^{-\sigma t} dt = \int_{-\infty}^{0} e^{-(\sigma + 1)t} dt + \int_{0}^{+\infty} e^{-(\sigma + 2)t} dt
    \]
    * برای همگرایی جمله اول در $t \to -\infty$ نیاز داریم به: $\sigma + 1 < 0 \implies \sigma < -1$
    \newline
    * برای همگرایی جمله دوم در $t \to +\infty$ نیاز داریم به: $\sigma + 2 > 0 \implies \sigma > -2$
    \newline
    با اشتراک این دو شرط، بازه همگرایی برابر است با: $-2 < \sigma < -1$
\end{enumerate}

 با توجه به این که در تبدیل لاپلاس $\text{Re}(s) = \sigma$ است، مشاهده می‌شود که بازه‌های مجاز $\sigma$ برای همگرایی انتگرال، دقیقاً و به طور کامل با نواحی همگرایی به دست آمده در بخش الف مطابقت دارند.


\subsectionaddtolist{ج}

تبدیل فوریه یک سیگنال پدید می‌آید اگر بتوان در تبدیل لاپلاس آن، مقدار $\sigma = 0$ (یا همان $s = j\omega$) را قرار داد. این امر بدین معناست که محور موهومی ($s = j\omega$) باید درون ناحیه همگرایی سیستم قرار گرفته باشد.

با بررسی سه حالت ROC به دست آمده:
\begin{itemize}
    \item $\sigma > -1$: مقدار $\sigma = 0$ در این بازه قرار دارد. بنابراین برای سیگنال $x_1(t)$ {تبدیل فوریه وجود دارد}.
    \item $\sigma < -2$: مقدار $\sigma = 0$ در این بازه قرار ندارد. بنابراین برای سیگنال $x_2(t)$ {تبدیل فوریه وجود ندارد}.
    \item $-2 < \sigma < -1$: مقدار $\sigma = 0$ در این بازه قرار ندارد. بنابراین برای سیگنال $x_3(t)$ {تبدیل فوریه وجود ندارد}.
\end{itemize}

حضور محور موهومی ($\sigma = 0$) در ناحیه همگرایی، معادل با برقراری شرط {انتگرال‌پذیری مطلق} سیگنال در حوزه زمان است، یعنی:
\[
\int_{-\infty}^{+\infty} |x(t)| dt < \infty
\]
این شرط فیزیکی تضمین می‌کند که سیگنال پایدار بوده و انرژی یا رفتار آن در بی‌نهایت محدود و میرا شونده است. اگر محور موهومی در ROC نباشد (مثل حالت ۲ و ۳)، سیگنال در زمان‌های مثبت یا منفی به صورت نمایی رشد کرده و واگرا می‌شود؛ در نتیجه مفهوم فیزیکی تبدیل فوریه (که بر پایه پایداری و تحلیل فرکانسی سیگنال‌های با انرژی یا توان محدود است) برای آن برقرار نخواهد بود.
